\documentclass[11pt,a4paper]{article}
\usepackage{fontspec}
\usepackage[margin=1in]{geometry}
\usepackage[unicode,hidelinks]{hyperref}
\setlength{\parindent}{0pt}
\setlength{\parskip}{.6em}
\emergencystretch=3em
\newcommand{\volumeentry}[3]{\par\begingroup\parskip=.7em
  \noindent\href{Foundation_Trilogy_Complete_Collection.pdf\#page=#3}{\textbf{#1}\quad #2\nobreak\dotfill\nobreak\mbox{#3}}\par\endgroup}
\hypersetup{pdftitle={Foundation Trilogy},pdfauthor={Sidong Liu}}
% COVER SUBTITLE: Selection, sharing and physical return
\begin{document}
\begin{titlepage}
\centering
\vspace*{1.6cm}
{\Huge\bfseries Foundation Trilogy\par}
\vspace{1cm}
{\Large Selection, sharing and physical return\par}
\vfill
{\Large Complete Collected Volume\par}
\vspace{1.6cm}
{\Large Sidong Liu, PhD\par}
\vspace{.3cm}
{iBioStratix Ltd\par}
\href{mailto:sidongliu@hotmail.com}{\texttt{sidongliu@hotmail.com}}\par
\vspace{1cm}
{\large October 2026\par}
\end{titlepage}
\setcounter{page}{2}
\section*{Contents}
The numbers below refer to pages in this PDF. Each paper retains its own printed pagination, proofs, appendices and references.
\volumeentry{}{Introduction and reading map}{3}
\volumeentry{F-I}{Foundation I: From Source to Manifestation. Accessible Algebras from a Source, Its Noise and a Readable Seed}{4}
\volumeentry{F-II}{Foundation II: Shared Manifestation. Shared Structure, Public Calibration and Different Histories}{30}
\volumeentry{F-III}{Foundation III: The Perfuming Loop. Held Content, Learned Frames and Physical Return}{47}
\clearpage
\section*{Introduction}\label{collection:intro}
Foundation asks whether the structures used elsewhere in the theory can exist, be shared and return into further physical action. A named algebra is not enough: its source, permitted operations, reference state and representation determine which conclusions can actually be inherited.

I connects source and noise to accessible hulls. Its finite centre example and conformal construction have distinct types; an internal standard Cut requires the latter's representation and modular conditions. II derives common structure under stated symmetry assumptions, separates agreement from correctness, and constructs a common finite-energy pulse whose weak-mixture tangent can enter the source-identification test of Consciousness II. Positive energy fixes a time orientation, not every charge orientation.

III connects held content, history-selected protection and physical return. In one finite semiclassical mean-action scalar/JT realization, the acquired record changes the apparatus source and a later X-basis probe receives the response in the same history. Its approximately $10^{-9}$ signal is an existence witness, not a useful amplifier. Separate quantum-material renewal constructions remain separately accounted for. The finite A result neither requires nor proves fixed-capacity sustained thermal updating B or joint sharp heat export and effective return C.

\section*{Reading map}
\begin{description}
\item[F-I] Source, noise and authorized seeds generate hulls of different types. Geometric inheritance requires the declared state, representation and internal standard Cut.
\item[F-II] Shared structure, shared state and correct public content are distinguished. Common finite-energy preparations provide a physical source-identification interface.
\item[F-III] Held content and finite physical return are realized with explicit resources. A history-selected frame and protected old quantum content are not the same code.
\end{description}
\end{document}
